\subsection{2026-09-22 --- Multi-horizon features for the count model}
\label{sec:countfeat}

The count is the binding constraint (Sections~\ref{sec:oracle}--\ref{sec:robustness}). This entry
targets it directly with a multi-horizon time-series feature set, built as
\texttt{src/count\_features.py} and evaluated at \textbf{(case, unit)} granularity on
\textbf{2021} --- the test-like year, following Section~\ref{sec:robustness}.

\subsubsection*{Design.}
\begin{itemize}\itemsep1pt
  \item \textbf{Windows in both directions}: trailing 1w/1m/3m/1y, forward 1w/1m, and previous and
        next week explicitly. Water is storable, so SHOP trades volume across weeks --- high prices
        next week should suppress this week's count. Forward windows are permitted; only
        test-period targets are withheld (Section~\ref{sec:eda-round-3b}).
  \item \textbf{Relative rather than absolute}: for every horizon, the week's level minus that
        horizon and the ratio to it, plus a seasonal comparison against the same calendar week
        across history. The spread changed sign after 2020, so absolute levels mislead across
        regimes (Section~\ref{sec:robustness}).
  \item \textbf{Inferred latent quantities}: SHOP sees demand and we do not, but price carries its
        signature --- rolling price volatility and seasonal price deviation stand in for demand
        pressure.
\end{itemize}

Applied to price, water value, spread, storage, inflow and price-volatility at system level (150
features), plus own-reservoir water value, own-spread and own-storage per plant (75 more). At
(case, unit) that is 244 features over 35{,}798 rows --- unlike the earlier count attempt, which
put 1{,}197 features over 2{,}557 cases (Section~\ref{sec:oracle}).

\subsubsection*{Result (validation 2021).}

\begin{center}
\begin{tabular}{lrrrr}
\toprule
Count features & Count MAE & 3-class acc & Blended score & $\Delta$ vs.\ pointwise \\
\midrule
Old (19)               & 48.8h & 0.616 & 0.8486 & $+0.0140$ \\
\textbf{New (244)}     & \textbf{43.2h} & \textbf{0.641} & \textbf{0.8561} & $\mathbf{+0.0216}$ \\
\bottomrule
\end{tabular}
\end{center}

\finding{\textbf{Count MAE falls 48.8h to 43.2h ($-11\%$) and it converts into score}: the
blended result rises 0.8486 to 0.8561, taking the budget augmentation from $+0.014$ to $+0.022$
over the 0.8346 pointwise baseline. This is the first feature work in several rounds that moved
the number on a test-like year.}

\finding{\textbf{The top count feature is an accidental water balance.}
\texttt{storage\_next\_minus\_prev} --- system storage after the case week minus storage before it
--- has gain 336{,}834, nearly twice the next feature. Net weekly drawdown \emph{is} water
released, which is production, which is count $\times$ capacity. The cascade water balance
hypothesis (Sections~\ref{sec:pruning}, \ref{sec:state}) therefore arrives here by a different
route, in its crudest possible form --- system-level, no cascade attribution, no upstream plant
discharge --- and still dominates every other count feature.}

\finding{\textbf{Gain share validates two design choices.} Relative-to-horizon comparisons are the
largest family at \textbf{33.1\%}, and own-reservoir features second at \textbf{27.7\%} --- the
per-unit specificity that only became expressible once \emph{unit} became a row dimension
(Section~\ref{sec:budget}). Forward windows read 6.9\% by family label, but that undercounts them:
\texttt{ownwv\_next\_week} is the third most important feature individually and the top feature
overall is built from next-week storage.}

\subsubsection*{Next.} The system-level drawdown proxy dominates while ignoring cascade structure
entirely. A per-plant balance --- natural inflow plus upstream plant discharge minus downstream
draw and spill, attributed through the topology graph --- is the obvious successor and is the last
major unexploited data source. Count MAE remains 43.2h against an 87h mean, and the oracle ceiling
is 0.9765, so the headroom is still largely intact.
